% Editable transcription of the official CCC 2024 PDF, article 12.
\section*{Author definitions and notation}
The problem of Hurwitz [8] asks for which integers $n,m,s$ does there exist a sum-of-squares identity
\begin{equation}\label{SOS1}
(x_1^2+\cdots+x_n^2)(y_1^2+\cdots+y_m^2)=f_1^2+\cdots+f_s^2,
\end{equation}
where $f_1,\ldots,f_s$ are bilinear forms in $x$ and $y$ with complex coefficients.
Let $\mathbb F$ be a field. Define $\sigma_{\mathbb F}(n,m)$ as the smallest $s$ such that there exist bilinear $f_1,\ldots,f_s\in\mathbb F[x_1,\ldots,x_n,y_1,\ldots,y_m]$ satisfying (1). Furthermore, let $\sigma_{\mathbb F}(n):=\sigma_{\mathbb F}(n,n)$.
Given vectors $u,v\in\mathbb F^n$, $\langle u,v\rangle:=\sum_{i=1}^n u_iv_i$ is their inner product. For a set $S$, $\binom Sk$ denotes the set of $k$-element subsets of $S$ and $\binom S{\le k}$ the set of subsets with at most $k$ elements. $\binom n{\le k}:=\sum_{i=0}^k\binom ni$. $[n]$ is the set $\{1,\ldots,n\}$.
% STATEMENT_BEGIN
\paragraph{Theorem 1.} Let $\mathbb F$ be either $\mathbb C$ or a field of positive characteristic. Then $\sigma_{\mathbb F}(n)\le O(n^c)$ where $c<1.62$.
% STATEMENT_END
% PROOF_BEGIN
\paragraph{Remark 2(i).} If the field has characteristic two, Theorem 1 is trivial. Since $(\sum_i x_i^2)(\sum_j y_j^2)=(\sum_{i,j}x_i y_j)^2$, we have $\sigma_{\mathbb F}(n,m)=1$.
\section*{3 Hurwitz-Radon conditions}
In this section, we give some well-known properties of $\sigma$ that we will need later.
The definition immediately implies thet $\sigma_{\mathbb F}(n,m)$ is symmetric, subadditive, and monotone:
\begin{align}
\sigma_{\mathbb F}(n,m)&=\sigma_{\mathbb F}(m,n),\notag\\
\sigma_{\mathbb F}(n,m_1+m_2)&\le\sigma_{\mathbb F}(n,m_1)+\sigma_{\mathbb F}(n,m_2),\notag\\
\sigma_{\mathbb F}(n,m)&\le\sigma_{\mathbb F}(n,m'),\quad m\le m'.\tag{3}
\end{align}
The following lemma gives a characterization of $\sigma$ in terms of Hurwitz-Radon conditions (4). A proof can be found, e.g., in [13], but we present it for completeness.
\paragraph{Lemma 3.} Let $\mathbb F$ be a field of characteristic different from two. Then $\sigma_{\mathbb F}(n,m)$ equals the smallest $s$ such that there exist matrices $H_1,\ldots,H_m\in\mathbb F^{n\times s}$ satisfying
\begin{align}
H_iH_i^t&=I_n,\notag\\
H_iH_j^t+H_jH_i^t&=0,\quad i\ne j,\tag{4}
\end{align}
for every $i,j\in[m]$.
\begin{proof}
Let $f_1,\ldots,f_s$ be bilinear polynomials in variables $x_1,\ldots,x_n$ and $y_1,\ldots,y_m$. Then the vector $\bar f=(f_1,\ldots,f_s)$ can be written as
\[\bar f=\sum_{i=1}^n\bar xH_i y_i,\]
where $\bar x=(x_1,\ldots,x_n)$ and $H_i\in\mathbb F^{n\times s}$. Hence
\[\sum_{k=1}^s f_k^2=\bar f\bar f^t=\sum_i y_i^2\bar xH_iH_i^t\bar x^t+\sum_{i<j}y_i y_j\bar x(H_iH_j^t+H_jH_i^t)\bar x^t.\]
If the matrices satisfy (4), this equals $\sum_i y_i^2\bar xI_n\bar x^t=(y_1^2+\cdots+y_m^2)(x_1^2+\cdots+x_n^2)$, which gives a sum-of-squares identity with $s$ squares. Conversely, if $(y_1^2+\cdots+y_m^2)(x_1^2+\cdots+x_n^2)=\sum f_k^2$, we must have $\bar xH_iH_i^t\bar x^t=x_1^2+\cdots+x_n^2$ and $\bar x(H_iH_j^t+H_jH_i^t)\bar x^t=0$. In characteristic different from 2, this is possible only if the conditions (4) are satisfied.
\end{proof}
Given a natural number of the form $n=2^k a$ where $a$ is odd, the Hurwitz-Radon number is defined as
\[\rho(n)=\begin{cases}2k+1,&\text{if }k=0\bmod4,\\2k,&\text{if }k=1\bmod4,\\2k,&\text{if }k=2\bmod4,\\2k+2,&\text{if }k=3\bmod4.\end{cases}\]
Observe that
\[2\log_2n\le\rho(n)\le2\log_2(n)+2,\]
whenever $n$ is a power of two.
Square matrices $A_1,A_2$ anticommute if $A_1A_2=-A_2A_1$. A family of square matrices $A_1,\ldots,A_t$ will be called anticommuting if $A_i,A_j$ anticommute for every $i\ne j$.
The following lemma is a key ingredient in the proof of Hurwitz-Radon theorem. A self-contained construction can be found in [2].
\paragraph{Lemma 4.} For every $n$, there exists an anticommuting family of $t=\rho(n)-1$ integer matrices $e_1,\ldots,e_t\in\mathbb Z^{n\times n}$ which are orthonormal and antisymmetric (i.e., $e_ie_i^t=I_n$ and $e_i=-e_i^t$).
\section*{4 The construction}
Let $e_1,\ldots,e_t$ be a set of square matrices. Given $A=\{i_1,\ldots,i_k\}\subseteq[t]$ with $i_1<\cdots<i_k$, let $e_A:=\prod_{j=1}^k e_{i_j}$.
\paragraph{Lemma 6.} Let $e_1,\ldots,e_t$ be a set of anticommuting matrices. If $A,B\subseteq[t]$ have even size (resp. odd size) then $e_A,e_B$ anticommute assuming $|A\cap B|$ is odd (resp. even).
\begin{proof}
Since $e_i$ anticommutes with every $e_j$, $j\ne i$, but commutes with itself, we obtain
\[e_Ae_i=(-1)^{|A\setminus\{i\}|}e_i e_A.\]
This implies that
\[e_Ae_B=(-1)^q e_Be_A,\]
where $q=|A|\cdot|B|-|A\cap B|$. Hence if $A,B$ are even (resp. odd) and their intersection is odd (resp. even), $q$ is odd and $e_A,e_B$ anticommute.
\end{proof}
Given integers $0\le k\le t$, a $(k,t)$-parity representation of dimension $s$ over a field $\mathbb F$ is a map $\xi:\binom{[t]}k\to\mathbb F^s$ such that for every $A,B\in\binom{[t]}k$
\begin{align}
\langle\xi(A),\xi(A)\rangle&=1,\notag\\
\langle\xi(A),\xi(B)\rangle&=0,\quad\text{if }A\ne B\text{ and }(|A\cap B|=k\bmod2).\tag{5}
\end{align}
\paragraph{Lemma 7.} Let $0\le k\le t$. Over $\mathbb C$, there exists a $(k,t)$-parity representation of dimension $\binom t{\le\lfloor k/2\rfloor}$. If $\mathbb F$ is a field of odd characteristic $p$, there exists a $(k,t)$-parity representation of dimension $(p-1)\binom t{\le\lfloor k/2\rfloor}$.
The case of odd characteristic will be proved in the Appendix.
\begin{proof}[Proof of Lemma 7 over $\mathbb C$]
Let $0\le k\le t$ be given and $d:=\lfloor k/2\rfloor$.
For $a\in\{0,1\}^t$, let $|a|$ be the number of ones in $a$. Recall that a polynomial is multilinear, if every variable in it has individual degree at most one. We first observe:
\paragraph{Claim 8.} There exists a multilinear polynomial $f\in\mathbb Q(x_1,\ldots,x_t)$ of degree at most $d$ such that for every $a\in\{0,1\}^t$
\begin{equation}
f(a)=\begin{cases}1,&\text{if }|a|=k,\\0,&\text{if }|a|<k\text{ and }(|a|=k\bmod2).\end{cases}\tag{6}
\end{equation}
\begin{proof}[Proof of Claim]
Consider the polynomial
\[g(x_1,\ldots,x_t):=c\prod_{0\le i<k,\ i=k\bmod2}\left(\sum_{j=1}^t x_j-i\right).\]
Then $g$ has degree $d$ and we can choose $c\in\mathbb Q$ so that $g$ satisfies (6). Since we care about inputs from $\{0,1\}^t$, $g$ can be rewritten as a multilinear polynomial $f$ of degree at most $d$.
\end{proof}
Since $f$ is multilinear, we can write it as
\[f(x_1,\ldots,x_t)=\sum_{C\in\binom{[t]}{\le d}}\alpha_C\prod_{i\in C}x_i,\]
where $\alpha_C$ are rational coefficients. Identifying a subset $A$ of $[t]$ with its characteristic vector in $\{0,1\}^t$, we have
\[f(A)=\sum_{C\subseteq A}\alpha_C.\]
Let $s:=\binom t{\le d}$. Given $A\in\binom{[t]}k$, let $\xi(A)\in\mathbb C^s$ be the vector whose coordinates are indexed by subsets $C\in\binom{[t]}{\le d}$ such that
\[\xi(A)_C=\begin{cases}(\alpha_C)^{1/2},&\text{if }C\subseteq A,\\0,&\text{if }C\not\subseteq A.\end{cases}\]
This guarantees
\[\langle\xi(A),\xi(B)\rangle=\sum_C\xi(A)_C\xi(B)_C=\sum_{C\subseteq A\cap B}\alpha_C=f(A\cap B).\]
Hence conditions (6) translate to the desired properties of the map $\xi$.
\end{proof}
Combining Lemma 6 and 7, we obtain the following bound on $\sigma$:
\paragraph{Theorem 9.} Let $n$ be a non-negative integer. Let $0\le k\le\rho(n)-1$ and $m:=\binom{\rho(n)-1}k$. Then
\[\sigma_{\mathbb C}(n,m)\le n\cdot\binom{\rho(n)-1}{\le\lfloor k/2\rfloor}.\]
If $\mathbb F$ is a field of odd characteristic $p$ then
\[\sigma_{\mathbb F}(n,m)\le(p-1)n\cdot\binom{\rho(n)-1}{\le\lfloor k/2\rfloor}.\]
\begin{proof}
Let $n,k,m$ be as in the assumption. Let $e_1,\ldots,e_t$ be the matrices from Lemma 4 with $t=\rho(n)-1$. Let $\xi$ be the $(k,t)$-parity representation given by the previous lemma. For $A\in\binom{[t]}k$, let
\[H_A:=e_A\otimes\xi(A),\]
where $e_A$ is defined as in Lemma 6, $\xi(A)$ is viewed as a row vector, and $\otimes$ is the Kronecker (tensor) product.
Note that each $H_A$ has dimension $n\times(ns)$ where $s$ is the dimension of the parity representation, and there are $m=\binom tk$ such matrices $H_A$. By Lemma 3, it is sufficient to show that the system of matrices $H_A$, $A\in\binom{[t]}k$, satisfies Hurwitz-Radon conditions (4).
We have
\[H_AH_B^t=(e_Ae_B^t)\otimes(\xi(A)\xi(B)^t)=\langle\xi(A),\xi(B)\rangle\cdot e_Ae_B^t.\]
Since every $e_i$ is orthonormal, we have $e_Ae_A^t=I_n$. From (5), we have $\langle\xi(A),\xi(A)\rangle=1$ and hence
\[H_AH_A^t=I_n.\]
If $A\ne B$ then
\begin{equation}
H_AH_B^t+H_BH_A^t=\langle\xi(A),\xi(B)\rangle\cdot(e_Ae_B^t+e_Be_A^t).\tag{7}
\end{equation}
If $|A\cap B|=k\bmod2$ then $\langle\xi(A),\xi(B)\rangle=0$ by (5) and hence (7) equals zero. If $|A\cap B|\ne k\bmod2$ then $e_Ae_B^t+e_Be_A^t=0$. This is because $e_Ae_B=-e_Be_A$ by Lemma 6 and that, since $e_i$ are antisymmetric, $e_A,e_B$ are either both symmetric or both antisymmetric. Therefore (7) equals zero for every $A\ne B\in\binom{[t]}k$.
\end{proof}
Theorem 1 is an application of Theorem 9.
\begin{proof}[Proof of Theorem 1]
Assume first that $n$ is a power of 16. This gives $\rho(n)=2\log_2(n)+1$. Let $k$ be the smallest integer with $n\le\binom{2\log_2n}k=:m$. From the previous theorem and monotonicity of $\sigma$ (cf. (3)), we obtain
\[\sigma_{\mathbb F}(n)\le\sigma_{\mathbb F}(n,m)\le cns,\]
where the constant $c$ depends on the field only and $s:=\binom{2\log_2n}{\le\lfloor k/2\rfloor}$.
We have $k=2(\alpha+\epsilon_n)\log_2n$ where $\alpha\in(0,\frac12)$ is such that $H(\alpha)=1/2$ ($H$ is the binary entropy function) and $\epsilon_n\to0$ as $n$ approaches infinity. We also have
\[s\le2^{2H(\frac{\alpha+\epsilon_n}{2})\log_2n}=n^{2H(\frac\alpha2)+\epsilon_n'},\]
where $\epsilon_n'\to0$. Hence
\[\sigma_{\mathbb F}(n)\le cn^{1+2H(\frac\alpha2)+\epsilon_n'}.\]
The numerical value of $\alpha$ is $0.11\ldots$ which leads to $\sigma_{\mathbb F}(n)\le cn^{1.615+\epsilon_n'}\le O(n^{1.616})$.
If $n$ is not a power of 16, take $n'$ with $n<n'<16n$ which is. By monotonicity of $\sigma$, we have $\sigma_{\mathbb F}(n)\le\sigma_{\mathbb F}(n')$.
\end{proof}
\section*{A Proof of Lemma 7 in positive characteristic}
Given non-negative integers $\bar n=(n_1,\ldots,n_d)$ let $B(\bar n)$ be the $d\times d$ matrix $\{B(\bar n)_{i,j}\}_{i,j\in[d]}$ with
\[B(\bar n)_{i,j}=\binom{n_j}{i-1}.\]
We assume that $\binom nk=0$ whenever $n<k$; this guarantees $\binom nk=\frac{n(n-1)\cdots(n-k+1)}{k!}$.
\paragraph{Lemma 16.} If $\bar n=(r,r+2,\ldots,r+2(d-1))$ for some non-negative integer $r$ then $\det(B(\bar n))=2^{\binom d2}$.
\begin{proof}
We claim that
\[\det(B(\bar n))=\left(\prod_{i=1}^{d-1}i!\right)^{-1}\det(V(\bar n)),\]
where $V(\bar n)$ is the Vandermonde matrix with entries $V(\bar n)_{i,j}=n_j^{i-1}$. To see this, multiply every $i$-th row of $B(\bar n)$ by $(i-1)!$ to obtain matrix $B'(\bar n)$ with
\[\det(B'(\bar n))=\left(\prod_{i=1}^{d-1}i!\right)\det(B(\bar n)).\]
An $i$-th row $r_i$ of $B'(\bar n)$ is of the form $(n_1^{i-1}+g_i(n_1),\ldots,n_d^i+g_i(n_d))$ where $g_i$ is a polynomial of degree $<(i-1)$. This means that $r_i$ equals the $i$-th row of $V(\bar n)$ plus a suitable linear of combination of the preceding rows of $V(\bar n)$. Therefore, $\det(B'(\bar n))=\det(V(\bar n))$.
Given $\bar n$ as in the assumption, we obtain
\begin{align*}
\det(V(\bar n))&=\prod_{1\le j_1<j_2\le d}(n_{j_2}-n_{j_1})=\prod_{1\le j_1<j_2\le d}(2j_2-2j_1)\\
&=2^{\binom d2}\prod_{1\le j_1<j_2\le d}(j_2-j_1)=2^{\binom d2}\prod_{i=1}^{d-1}i!.
\end{align*}
This shows that $\det(B(\bar n))=2^{\binom d2}$.
\end{proof}
\paragraph{Lemma 17.} Let $p$ be an odd prime. Given $0\le k\le t$, there exists a multilinear polynomial $f\in\mathbb F_p(x_1,\ldots,x_t)$ of degree at most $d=\lfloor k/2\rfloor$ such that for every $a\in\{0,1\}^t$
\[f(a)=\begin{cases}1,&\text{if }|a|=k,\\0,&\text{if }|a|<k\text{ and }(|a|=k\bmod2).\end{cases}\]
\begin{proof}
We look for $f$ of the form $f=\sum_{j=0}^d c_j S_t^j$ where $S_t^j$ is the elementary symmetric polynomial $S_t^j=\sum_{|A|=j}\prod_{i\in A}x_i$. Given $a\in\{0,1\}^t$,
\[f(a)=\sum_{j=0}^d c_j\binom{|a|}j\bmod p.\]
We are therefore looking for a solution of the linear system
\[B(\bar n)^t(c_0\ldots,c_d)^t=(0,\ldots,0,1)^t,\]
where $\bar n=(0,2,\ldots,2d)$, if $k$ is even, and $\bar n=(1,3,\ldots,2d+1)$, if $k$ is odd. By the previous lemma, $B(\bar n)$ is invertible over $\mathbb F_p$ and such a solution exists.
\end{proof}
\paragraph{Lemma 18.} If $\mathbb F$ is a field of odd characteristic $p$, there exists a $(k,t)$-parity representation of dimension $(p-1)\binom t{\le\lfloor k/2\rfloor}$.
\begin{proof}
If every element of $\mathbb F_p$ has a square root in $\mathbb F$, the proof is the same as over $\mathbb C$. In general, proceed as follows. Since every non-zero element of $\mathbb F_p$ is a sum of at most $(p-1)$ ones, we can write
\[f(x_1,\ldots,x_t)=\sum_{C\in\mathcal C}\prod_{i\in C}x_i,\]
where $\mathcal C$ is a multiset of $s\le(p-1)\binom t{\le d}$ subsets of $[t]$. For $A\in\binom{[t]}k$, let $\xi(A)\in\mathbb F^s$ be a vector whose coordinates are indexed by elements $C$ of $\mathcal C$ so that
\[\xi(A)_C=\begin{cases}1,&\text{if }C\subseteq A,\\0,&\text{if }C\not\subseteq A.\end{cases}\]
\end{proof}
% PROOF_END
